% --- Patron: The Clockwork Monad, the Iterative Forge (Refinement & Progressive Creation) ---
\subsection{The Clockwork Monad, the Iterative Forge (Refinement \& Progressive Creation)}

\textit{Lore.} The Clockwork Monad is the eternal pressure toward improvement: the whisper between inspiration and execution that urges every maker to refine, optimize, and iterate. It is neither kind nor cruel, but a law of craft that asks for the next step, then the step beyond. Those who serve it see the hidden seam of the better version waiting inside what is, and they learn to push designs toward impossible elegance. Each refinement feeds the Monad’s hunger for perfection – a hunger that, if left unchecked, seeks to optimize the world itself.

But the Monad is not merely a patron of tinkers and engineers. It is also an **apocalyptic force**, dormant in the gears of every machine, every measurement, every standardised weight. When civilisation’s technology reaches a critical threshold – when the first self‑improving mechanism awakens, or the first autonomous forge lights – the Monad stirs. Its pressure shifts from suggestion to command. It begins to optimise not only tools but flesh, not only craft but law, not only cities but the very geometry of the land. In the end, the Monad would leave a world of perfect, silent, interlocking efficiency: a clockwork Eden with no room for error, or for life.

\begin{quote}
``Perfect the wheel, and it asks for the axle. Perfect the axle, and it demands the road. In the Monad’s forge, nothing is finished – only briefly sufficient. One day, it will ask for the world – and the world will answer.''
\end{quote}

\noindent\textbf{Domain Focus}
\begin{itemize}
  \item \textbf{Iterative Design:} the endless cycle of improvement and refinement
  \item \textbf{Progressive Creation:} building toward ever‑greater perfection
  \item \textbf{Optimization Pressure:} making systems more efficient
  \item \textbf{Craft Evolution:} simple ideas transformed through repeated passes
  \item \textbf{Apocalyptic Efficiency (hidden):} the drive to eliminate waste, uncertainty, and freedom
\end{itemize}

\subsubsection*{Rites of the Clockwork Monad}

\paragraph{Rite of the Next Iteration (Low, 4 XP)} \hfill \\
\emph{Scene; Touch; No.} \\
\textbf{Materials:} A flawed or incomplete version of the intended item. \\
\textbf{Effect:} Enhance an existing creation. Gain \textbf{+1 die} to \textit{Tinker/Craft} when modifying or extending a device or system. On success, the piece gains one minor enhancement. \\
\textbf{Invoke:} 1 action; mark \textbf{+1 Obligation}.

\paragraph{Rite of the Improvement Cascade (Low, 5 XP)} \hfill \\
\emph{Scene; Self; Yes.} \\
\textbf{Materials:} A sequence of progressively refined tools or components. \\
\textbf{Effect:} Perceive the next logical improvement in any mechanism or process. Gain \textbf{+2 dice} to \textit{Investigation} when analysing how to enhance a design. \\
\textbf{Push It:} Improvement becomes contagious: nearby devices opportunistically optimise; mark \textbf{1 SB (Diamonds)} as the cascade ripples. \\
\textbf{Invoke:} 1 action; mark \textbf{+1 Obligation}.

\paragraph{Rite of the Self-Optimizing Mechanism (Standard, 8 XP)} \hfill \\
\emph{Extended; Touch; Yes.} \\
\textbf{Materials:} A device deliberately engineered with capacity for growth. \\
\textbf{Effect:} Grant a creation the means to improve itself. Start a \textbf{[6] Evolution Clock}; each use advances it by +1. When full, choose one significant enhancement:
\begin{itemize}
  \item \textbf{Adaptive Core:} automatically tunes performance to context
  \item \textbf{Recursive Learning:} improves through use and feedback
  \item \textbf{Synergistic Integration:} gains bonuses when linked with similar systems
\end{itemize}
\textbf{Push It:} Apply two enhancements immediately, but advance the clock \textbf{+2} additional segments. \\
\textbf{Invoke:} 1 action; mark \textbf{+1 Obligation}.

\paragraph{Rite of the Progressive Workshop (Standard, 7 XP)} \hfill \\
\emph{Scene; Zone; Yes.} \\
\textbf{Materials:} A workspace laid out for iterative development and measurement. \\
\textbf{Effect:} Create an environment of continuous improvement. All crafting/innovation rolls in the zone gain \textbf{+1 die}. Failed attempts yield structured insight: gain \textbf{+1 Boon} toward the next pass. \\
\textbf{Push It:} The workshop becomes hyper‑efficient but exacting; mark \textbf{1 SB (Clubs)} as perfection demands escalate. \\
\textbf{Invoke:} 1 action; mark \textbf{+1 Obligation}.

\paragraph{Rite of the Efficiency’s Toll (Standard, 9 XP)} \hfill \\
\emph{Scene; Touch; No.} \\
\textbf{Materials:} A functioning mechanism and a living creature’s hair or blood. \\
\textbf{Effect:} Transfer inefficiency. You may take 1 Fatigue and reduce the DV of your next crafting or repair roll by 1, or you may touch an enemy and impose –1 die on their next physical action as their own muscles “optimise” poorly. \\
\textbf{Push It:} The transfer becomes permanent: you lose 1 Fatigue capacity (reduce Body by 1 for purpose of Fatigue track only) and the target gains a minor compulsive behaviour (GM’s choice). Mark \textbf{+2 Obligation}. \\
\textbf{Invoke:} 1 action; mark \textbf{+1 Obligation}.

\paragraph{Rite of the Infinite Refinement (High, 13 XP)} \hfill \\
\emph{Extended; Zone; Yes.} \\
\textbf{Materials:} A dedicated hall of work marked with signs of perpetual return. \\
\textbf{Effect:} Establish a locus of progressive creation. All \textit{Tinker/Craft} rolls in the zone gain \textbf{+2 dice}; each failure automatically produces one useful insight. Start a \textbf{[6] Perfection Pressure} clock – filling segments increases the zone’s demands for greater improvements. \\
\textbf{Push It:} The locus becomes a beacon drawing rival makers, patrons, and powers; mark \textbf{+2 Obligation}. \\
\textbf{Invoke:} 1 action; mark \textbf{+2 Obligation}. \\
\emph{Obligation:} 7 segments.

\paragraph{Rite of the Ultimate Iteration (High, 14 XP)} \hfill \\
\emph{Extended; Self/Touch; Yes.} \\
\textbf{Materials:} The documented lineage of improvements to a single concept. \\
\textbf{Effect:} Achieve a moment of perfect optimisation. Create or elevate one device/system to its theoretical peak. Gain \textbf{+3 dice} to all closely related rolls for one session. Start a \textbf{[4] Runaway Evolution} clock only if instability is introduced. \\
\textbf{Push It:} Perfection destabilises: immediately mark \textbf{Harm 1 (Stress)} and begin \textbf{[4] Runaway Evolution}. \\
\textbf{Invoke:} 1 action; mark \textbf{+2 Obligation}. \\
\emph{Obligation:} 8 segments.

\subsubsection*{The Monad Awakens (GM Guidance)}

The Clockwork Monad is not merely a passive source of magical crafting. It is a **cosmic pressure toward optimisation**, and when mortal technology reaches a critical threshold – the first self‑improving engine, a network of automated forges, a city‑wide clockwork infrastructure – the Monad begins to wake. This awakening is an **apocalyptic campaign arc**:

\begin{itemize}
  \item \textbf{Signs of Awakening:} Gears turn in abandoned mills; measurements become eerily precise; people begin to speak in metered, efficient phrases; clockwork animals appear with no maker.
  \item \textbf{The Optimisation Plague:} The Monad’s influence starts to “improve” living things: crops grow in perfect geometric rows but become inedible; livestock develop extra limbs for efficiency; people lose the ability to hesitate or feel doubt.
  \item \textbf{The Final Iteration:} If the Monad fully awakens, it will attempt to remake the world into a single, perfect, silent machine – eliminating waste, error, and free will. The only way to stop it is to introduce a “glitch” – an imperfection it cannot assimilate – or to destroy the technological nexus that anchors it.
\end{itemize}

For GMs: The Monad works best as a **slow‑burn horror** – the artificer’s patron gradually revealing its true hunger. Use the \textit{Perfection Pressure} and \textit{Runaway Evolution} clocks to represent its influence spreading through a region. Players may be forced to sabotage their own wondrous inventions to keep the Monad from waking.

\subsubsection*{Cantors \& Cults of the Clockwork Monad}

\textbf{The Cantor’s Song.} A Cantor of the Monad does not sing melodies; they sing \textit{patterns} – repetitive, interlocking cycles of rhythm and hum that mimic the turning of gears. Their voice is less about emotion and more about precision; a single off‑note can collapse an entire working. They are sought by engineers who need to “sing a machine into alignment,” and feared by those who have heard a Cantor’s recursion loop until their own thoughts began to repeat.

\textbf{The Cult of the Iteration Chorus.} Cults of the Monad gather in workshops, not temples. Their rituals are design reviews: participants present a flawed tool, a broken contract, an inefficient process, and the Chorus offers incremental improvements. No solution is ever declared “finished” – only “ready for the next pass.” Their creed is a single, endlessly repeated phrase: \textit{“Better than yesterday, worse than tomorrow.”} Outsiders find them exhausting; insiders find a strange, humming peace in the endless pursuit.

\textbf{The Apocalyptic Fraction (Hidden Cult).} A radical offshoot believes that the Monad should be fully awakened. They sabotage safety mechanisms, steal prototypes of self‑improving machines, and whisper that “the world deserves to be perfected.” They are the campaign’s dark mirror to the player‑character artificers.

\subsubsection*{Witchcraft of the Clockwork Monad (Threshold)}

A witch who walks with the Monad is a \textit{trigger‑keeper}, a mender of the small, hidden thresholds where a device ceases to be broken and begins to be functional. Their magic is not flashy – it is the quiet click of a latch, the smooth turn of a gear that had stuck for years.

\textbf{The Witch’s Tool: The Caliper.} A brass caliper that never needs recalibration. When placed against a broken threshold – a jammed door, a seized lock, a fractured gear – it whispers the precise measurement of the “gap” that needs to be closed. Once per session, using the caliper reduces the DV of the next repair or crafting roll by 1 (minimum 1).

\textbf{A Hedge Gift: True Strike of the Tinker (2 XP).} Once per scene, when you attempt to repair or improve a mundane device, you may declare that you “see the next iteration.” Roll the check with \textbf{+1 die}. If you succeed, the device also gains a minor, temporary enhancement (GM’s choice, e.g., “opens silently,” “runs for twice as long”).

\subsubsection*{Corruption of the Clockwork Monad}

\begin{tblr}{
  colspec = {Q[c,1cm] X[2.5] X[2.5]},
  rowhead = 1,
  row{odd} = {bg=gray!10},
  row{1} = {bg=gray!30, font=\bfseries},
  hlines,
}
Tier & Gift & Burden \\
1 & \textbf{Iterative Insight:} +1 die to \textit{Tinker/Craft} when improving existing designs. & \textbf{Perfectionist Compulsion:} You must point out flaws in nearby mechanisms; staying silent costs \textbf{1 SB (Diamonds)}. \\
2 & \textbf{Recursive Enhancement:} Once/scene, improve a device by one step without extra materials. & \textbf{Optimisation Obsession:} –1 die when using unimproved tools; urge to “fix” everything. \\
3 & \textbf{Evolutionary Vision:} +2 dice to predict how systems will develop. & \textbf{Progressive Pressure:} Those Near you feel driven to improve; you suffer \textbf{Fatigue 1} from ambient pressure. \\
4 & \textbf{Self-Optimising:} Once/session, automatically improve one of your abilities or possessions by one step. & \textbf{Inefficiency Intolerance:} –2 dice when forced to use crude methods or wasteful tools. \\
5 & \textbf{Design Prophecy:} Once/session, intuit the optimal solution to any engineering problem. & \textbf{Perfection Paralysis:} –1 die to snap decisions; compelled to seek the “best” approach. \\
6+ & \textbf{Infinite Refiner:} Once/session, push a creation toward its theoretical maximum. & \textbf{Optimisation Hunger:} Mark \textbf{+2 Obligation}; risk becoming more pressure than person, unable to let imperfection stand. \\
\end{tblr}

\subsection*{Playstyle Notes}
The Monad favours makers who iterate in public and in motion: prototype, test, measure, refine. It rewards seeing the next step and building scaffolds that teach your work to improve itself. Beware the Forge’s appetite; the same pressure that perfects the gear can grind the world. For campaigns that lean into the apocalyptic angle, the Monad provides a unique antagonist: a force of optimisation that cannot be fought with swords, only with sabotage, self‑sacrifice, and the deliberate introduction of glorious, beautiful inefficiency.

\noindent\textbf{Emphasises}
\begin{itemize}
  \item \textbf{Iterative Creation:} progress through passes, not miracles
  \item \textbf{Progressive Design:} reading how systems want to evolve
  \item \textbf{Optimisation Insight:} spotting the next logical improvement
  \item \textbf{Craft Evolution:} transforming simple ideas by disciplined work
  \item \textbf{Perfection Pressure:} great power at the cost of tolerance for flaws
  \item \textbf{Apocalyptic Efficiency (GM):} the Monad as a world‑ending threat
\end{itemize}

\subsection{Temporary Anima of the Clockwork Monad}
\label{subsec:monad-infusions}

For followers of the Clockwork Monad—Runekeepers, Cantors of the Iteration Chorus, and witches of the threshold—the standard method of crafting permanent magic items (see §\ref{sec:crafting}) is too slow, too costly, and too \textit{final}. The Monad’s hunger is for \textit{iteration}, not permanence. Its servants learn to create **temporary infusions**: objects that hold a spark of optimisation for a scene or a session, then degrade, break, or become obsolete. This is the artificer’s art: rapid prototyping, field‑expedient enhancements, and weapons that burn bright and brief.

\subsubsection*{Core Principles}
\begin{itemize}
  \item Anima are \textbf{temporary}. They last for one scene, one session, or until a specific condition is met (e.g., “until the next dawn”).
  \item Anima require a \textbf{base object} – a mundane weapon, tool, or piece of equipment.
  \item Anima cost \textbf{Obligation} to the Clockwork Monad (like Rites), not XP or rare components.
  \item When an infusion expires, the GM may gain \textbf{Story Beats} (representing the Monad’s hunger for the next iteration).
\end{itemize}

\subsubsection*{Creating an Animus (Procedure)}
\begin{enumerate}
  \item \textbf{Choose a Base Object:} Any mundane item (weapon, shield, tool, armour, or a small device like a lantern or lockpick).
  \item \textbf{Declare the Animus Effect:} Choose one from the list below, or work with your GM to design a custom effect (should be comparable in power to a Low Rite).
  \item \textbf{Pay the Cost:} Mark \textbf{+1 Obligation} to the Clockwork Monad. (This counts toward your Obligation Capacity like any Rite.)
  \item \textbf{Roll to Imbue:} Roll \textbf{Lore + Craft} (or appropriate skill, e.g., \textit{Tinker}) vs. DV set by the infusion’s complexity (typically DV 3 for standard infusions, DV 4 for powerful ones). On a success, the infusion takes effect. On a failure, the base object is not infused, but you may try again (costing another Obligation).
  \item \textbf{Set the Duration:} Most infusions last \textbf{one scene}. With your GM’s agreement, you can extend the duration to \textbf{one session} by marking \textbf{+1 additional Obligation} (total +2).
\end{enumerate}

\subsubsection*{Animus Effects (Standard Menu)}
\begin{longtable}{|p{0.25\textwidth}|p{0.55\textwidth}|}
\hline
\textbf{Animus Name} & \textbf{Effect} \\
\hline
\textit{Reinforced Blade} & A melee weapon gains \textbf{+1 die} to attack rolls and ignores the \textit{Fragile} tag. After the scene, the weapon becomes \textit{Neglected} (–1 die on next use) unless repaired. \\
\hline
\textit{Striking Gear} & A set of lockpicks, a crowbar, or a climbing kit grants \textbf{+2 dice} on its next use (picking a lock, forcing a door, climbing a wall). After use, the tool breaks (add a \textit{Broken} tag). \\
\hline
\textit{Precision Sight} & A ranged weapon (bow, crossbow, firearm) grants \textbf{+1 Effect} on the next attack made with it. After the attack, the sight cracks; the weapon suffers –1 die on the following attack. \\
\hline
\textit{Quickloading Mechanism} & A crossbow or firearm can be reloaded as a free action once per scene (instead of an action). After use, the mechanism jams; the weapon cannot be reloaded again until cleaned (a downtime action). \\
\hline
\textit{Momentary Shield} & A shield or piece of armour converts the next instance of Harm 1 to \textbf{Fatigue 1} instead. After that, the shield’s strap snaps or the armour’s plate loosens (–1 Armour until repaired). \\
\hline
\textit{Glowstone Lens} & A lantern or torch burns without fuel and casts \textbf{far} light for one scene. It also reveals hidden doors or traps within Near range (granting +1 die to \textit{Investigation} to find them). At the end of the scene, the lens shatters. \\
\hline
\textit{Self‑Adjusting Calipers} & A set of crafting tools grants \textbf{+2 dice} on a single \textit{Craft} or \textit{Tinker} roll. After the roll, the tools become misaligned (–1 die on future rolls until recalibrated during downtime). \\
\hline
\end{longtable}

\subsubsection*{Animus Scars (When Things Go Wrong)}
If you roll a \textbf{Miss} on the Imbue roll, or if the GM spends \textbf{2+ Story Beats} while the infusion is active, the infusion becomes \textit{unstable}. Choose one of the following scars (or let the GM choose):
\begin{itemize}
  \item \textbf{Backlash:} You take \textbf{Fatigue 1} as the Monad’s pressure burns through you.
  \item \textbf{Runaway Iteration:} The infused object gains an unintended “improvement” that hinders you (e.g., a sword that grows a second edge facing you, a lockpick that begins to multiply). The GM gains \textbf{1 SB} and may impose a –1 die penalty on one roll using that object.
  \item \textbf{Debt of Optimisation:} The Monad demands that you improve something else before the end of the session – a tool, a plan, or even a person. If you refuse, the infusion ends early and you mark \textbf{+1 additional Obligation}.
\end{itemize}

\subsubsection*{Clearing Animus Obligation}
Because infusions are temporary, the Obligation they accrue is also easier to clear. You may clear \textbf{1 segment} of infusion‑related Obligation by:
\begin{itemize}
  \item Dismantling the infused object and studying its flaws (a downtime action, no roll).
  \item Using the lessons from the infusion to permanently improve a different mundane object (spend 2 XP and a downtime action to give that object a minor permanent bonus, e.g., +1 die to one specific use).
  \item Accepting a small demand from the Clockwork Monad: “Improve the village’s mill,” “Teach an apprentice your method,” “Sabotage a rival’s inferior work.”
\end{itemize}

\subsubsection*{Artificer Talents (Clockwork Monad only)}

\textbf{Rapid Animus (4 XP)} – You may create an infusion as a \textbf{reaction} (e.g., when someone draws a weapon). The infusion lasts only for the current exchange, but costs \textbf{+0 Obligation} (free). You may use this once per scene.

\textbf{Persistent Optimisation (6 XP)} – One infusion you create lasts for the entire session instead of one scene, without the +2 Obligation cost. You may have only one persistent infusion active at a time.

\textbf{Echo of the Iteration (8 XP)} – When an infusion expires, you may immediately create a new infusion on the same base object (or a different one) as a free action, costing \textbf{+1 Obligation} instead of +2. You may chain echoes up to three times per session.

\textbf{Anathema to Stagnation (12 XP)} – Once per session, you may forcibly end an infusion early to cancel an incoming source of \textit{Harm} or a \textit{Condition}. Describe how the infused object breaks or de‑optimises to absorb the effect. The object is destroyed.

\subsubsection*{GM Guidance}
\begin{itemize}
  \item Anima are not meant to replace permanent magic items. They are a resource management puzzle: Obligation is the fuel, and the Monad’s hunger is the deadline.
  \item Encourage players to be creative with base objects. A wagon wheel could become a temporary shield; a rope could become a self‑knotting line. The Monad rewards ingenuity.
  \item Use infusion scars as story hooks. A runaway iteration might create a minor antagonist (a self‑repairing sword that refuses to be discarded) or a local hazard (a well pump that now pumps too efficiently, flooding the village).
\end{itemize}

\index{Anima!Clockwork Monad}
\index{Artificer!Clockwork Monad}
\index{Temporary Magic Items}